\documentclass[10pt]{article}
\usepackage[T1]{fontenc}
\usepackage[margin=1in]{geometry}
\usepackage{amsmath,amssymb,booktabs,xcolor,enumitem,textcomp}
\usepackage[hidelinks]{hyperref}
\newcommand{\du}{\Delta u_{\mathrm{c}}}
\newcommand{\deps}{\Delta\varepsilon_{\mathrm{in}}}
% GENERATED by scripts/tcpmt_manuscript_macros.py -- do not edit
% source: /home/beams0/SABDELZAHER/ETM_project/HVR_OPT_work/research/final_closure
\newcommand{\PlossZone}{144.3}
\newcommand{\EffZone}{87.4}
\newcommand{\PconvZone}{140.7}
\newcommand{\ProuteZone}{3.62}
\newcommand{\PhvZone}{1.77}
\newcommand{\PtsvZone}{0.00}
\newcommand{\PattZone}{--}
\newcommand{\PintraZone}{1.81}
\newcommand{\PboardZone}{--}
\newcommand{\PpadsZone}{0.044}
\newcommand{\TmaxZone}{115.7}
\newcommand{\TdieZone}{89.8}
\newcommand{\TconvZone}{107.7}
\newcommand{\TjointZone}{115.0}
\newcommand{\WarpZone}{25.5}
\newcommand{\PkgWarpZone}{25.2}
\newcommand{\DuZone}{2.21}
\newcommand{\EpsZone}{5.71}
\newcommand{\EpsUZone}{1.7}
\newcommand{\EtItersZone}{5}
\newcommand{\TjointHotZone}{115.4}
\newcommand{\TlsZone}{109.7}
\newcommand{\TlsRiseZone}{85}
\newcommand{\ThsZone}{112.7}
\newcommand{\TindZone}{100.7}
\newcommand{\QplateZone}{--}
\newcommand{\QtopZone}{--}
\newcommand{\QdieDownZone}{--}
\newcommand{\NjointZone}{1369}
\newcommand{\PrtZone}{94.1}
\newcommand{\UnderZone}{35}
\newcommand{\RiseZone}{50.2}
\newcommand{\PlossZoneSone}{179.9}
\newcommand{\EffZoneSone}{84.7}
\newcommand{\PconvZoneSone}{144.8}
\newcommand{\ProuteZoneSone}{35.13}
\newcommand{\PhvZoneSone}{31.48}
\newcommand{\PtsvZoneSone}{0.01}
\newcommand{\PattZoneSone}{--}
\newcommand{\PintraZoneSone}{3.60}
\newcommand{\PboardZoneSone}{--}
\newcommand{\PpadsZoneSone}{0.045}
\newcommand{\TmaxZoneSone}{140.5}
\newcommand{\TdieZoneSone}{91.9}
\newcommand{\TconvZoneSone}{115.4}
\newcommand{\TjointZoneSone}{140.3}
\newcommand{\WarpZoneSone}{32.3}
\newcommand{\PkgWarpZoneSone}{32.6}
\newcommand{\DuZoneSone}{2.72}
\newcommand{\EpsZoneSone}{8.14}
\newcommand{\EpsUZoneSone}{1.1}
\newcommand{\EtItersZoneSone}{5}
\newcommand{\TjointHotZoneSone}{140.3}
\newcommand{\TlsZoneSone}{116.7}
\newcommand{\TlsRiseZoneSone}{92}
\newcommand{\ThsZoneSone}{123.3}
\newcommand{\TindZoneSone}{105.2}
\newcommand{\QplateZoneSone}{--}
\newcommand{\QtopZoneSone}{--}
\newcommand{\QdieDownZoneSone}{--}
\newcommand{\NjointZoneSone}{1369}
\newcommand{\PrtZoneSone}{114.2}
\newcommand{\UnderZoneSone}{37}
\newcommand{\RiseZoneSone}{65.7}
\newcommand{\PlossZoneStwo}{296.6}
\newcommand{\EffZoneStwo}{77.1}
\newcommand{\PconvZoneStwo}{158.4}
\newcommand{\ProuteZoneStwo}{138.22}
\newcommand{\PhvZoneStwo}{130.28}
\newcommand{\PtsvZoneStwo}{0.02}
\newcommand{\PattZoneStwo}{--}
\newcommand{\PintraZoneStwo}{7.87}
\newcommand{\PboardZoneStwo}{--}
\newcommand{\PpadsZoneStwo}{0.046}
\newcommand{\TmaxZoneStwo}{225.1}
\newcommand{\TdieZoneStwo}{98.6}
\newcommand{\TconvZoneStwo}{140.7}
\newcommand{\TjointZoneStwo}{224.3}
\newcommand{\WarpZoneStwo}{54.7}
\newcommand{\PkgWarpZoneStwo}{57.0}
\newcommand{\DuZoneStwo}{4.41}
\newcommand{\EpsZoneStwo}{--}
\newcommand{\EpsUZoneStwo}{--}
\newcommand{\EtItersZoneStwo}{7}
\newcommand{\TjointHotZoneStwo}{224.3}
\newcommand{\TlsZoneStwo}{139.6}
\newcommand{\TlsRiseZoneStwo}{115}
\newcommand{\ThsZoneStwo}{158.3}
\newcommand{\TindZoneStwo}{120.0}
\newcommand{\QplateZoneStwo}{--}
\newcommand{\QtopZoneStwo}{--}
\newcommand{\QdieDownZoneStwo}{--}
\newcommand{\NjointZoneStwo}{1369}
\newcommand{\PrtZoneStwo}{162.5}
\newcommand{\UnderZoneStwo}{45}
\newcommand{\RiseZoneStwo}{134.1}
\newcommand{\PlossZtwo}{147.4}
\newcommand{\EffZtwo}{87.2}
\newcommand{\PconvZtwo}{144.8}
\newcommand{\ProuteZtwo}{2.67}
\newcommand{\PhvZtwo}{--}
\newcommand{\PtsvZtwo}{0.55}
\newcommand{\PattZtwo}{0.23}
\newcommand{\PintraZtwo}{1.84}
\newcommand{\PboardZtwo}{--}
\newcommand{\PpadsZtwo}{0.044}
\newcommand{\TmaxZtwo}{124.3}
\newcommand{\TdieZtwo}{90.0}
\newcommand{\TconvZtwo}{113.6}
\newcommand{\TjointZtwo}{97.2}
\newcommand{\WarpZtwo}{68.7}
\newcommand{\PkgWarpZtwo}{102.7}
\newcommand{\DuZtwo}{3.09}
\newcommand{\EpsZtwo}{7.85}
\newcommand{\EpsUZtwo}{0.8}
\newcommand{\EtItersZtwo}{5}
\newcommand{\TjointHotZtwo}{98.3}
\newcommand{\TlsZtwo}{116.7}
\newcommand{\TlsRiseZtwo}{92}
\newcommand{\ThsZtwo}{119.7}
\newcommand{\TindZtwo}{107.1}
\newcommand{\QplateZtwo}{--}
\newcommand{\QtopZtwo}{--}
\newcommand{\QdieDownZtwo}{--}
\newcommand{\NjointZtwo}{528}
\newcommand{\PrtZtwo}{93.5}
\newcommand{\UnderZtwo}{37}
\newcommand{\RiseZtwo}{54.0}
\newcommand{\PlossZtwoSone}{149.4}
\newcommand{\EffZtwoSone}{87.0}
\newcommand{\PconvZtwoSone}{145.0}
\newcommand{\ProuteZtwoSone}{4.42}
\newcommand{\PhvZtwoSone}{--}
\newcommand{\PtsvZtwoSone}{0.56}
\newcommand{\PattZtwoSone}{0.23}
\newcommand{\PintraZtwoSone}{3.59}
\newcommand{\PboardZtwoSone}{--}
\newcommand{\PpadsZtwoSone}{0.044}
\newcommand{\TmaxZtwoSone}{124.8}
\newcommand{\TdieZtwoSone}{90.1}
\newcommand{\TconvZtwoSone}{114.0}
\newcommand{\TjointZtwoSone}{97.4}
\newcommand{\WarpZtwoSone}{69.0}
\newcommand{\PkgWarpZtwoSone}{103.1}
\newcommand{\DuZtwoSone}{3.10}
\newcommand{\EpsZtwoSone}{7.88}
\newcommand{\EpsUZtwoSone}{0.8}
\newcommand{\EtItersZtwoSone}{5}
\newcommand{\TjointHotZtwoSone}{98.4}
\newcommand{\TlsZtwoSone}{117.2}
\newcommand{\TlsRiseZtwoSone}{92}
\newcommand{\ThsZtwoSone}{120.2}
\newcommand{\TindZtwoSone}{107.5}
\newcommand{\QplateZtwoSone}{--}
\newcommand{\QtopZtwoSone}{--}
\newcommand{\QdieDownZtwoSone}{--}
\newcommand{\NjointZtwoSone}{528}
\newcommand{\PrtZtwoSone}{94.8}
\newcommand{\UnderZtwoSone}{37}
\newcommand{\RiseZtwoSone}{54.7}
\newcommand{\PlossZtwoStwo}{153.7}
\newcommand{\EffZtwoStwo}{86.7}
\newcommand{\PconvZtwoStwo}{145.5}
\newcommand{\ProuteZtwoStwo}{8.16}
\newcommand{\PhvZtwoStwo}{--}
\newcommand{\PtsvZtwoStwo}{0.56}
\newcommand{\PattZtwoStwo}{0.23}
\newcommand{\PintraZtwoStwo}{7.33}
\newcommand{\PboardZtwoStwo}{--}
\newcommand{\PpadsZtwoStwo}{0.044}
\newcommand{\TmaxZtwoStwo}{125.9}
\newcommand{\TdieZtwoStwo}{90.3}
\newcommand{\TconvZtwoStwo}{114.9}
\newcommand{\TjointZtwoStwo}{97.7}
\newcommand{\WarpZtwoStwo}{69.7}
\newcommand{\PkgWarpZtwoStwo}{104.0}
\newcommand{\DuZtwoStwo}{3.11}
\newcommand{\EpsZtwoStwo}{7.95}
\newcommand{\EpsUZtwoStwo}{0.8}
\newcommand{\EtItersZtwoStwo}{5}
\newcommand{\TjointHotZtwoStwo}{98.8}
\newcommand{\TlsZtwoStwo}{118.1}
\newcommand{\TlsRiseZtwoStwo}{93}
\newcommand{\ThsZtwoStwo}{121.2}
\newcommand{\TindZtwoStwo}{108.1}
\newcommand{\QplateZtwoStwo}{--}
\newcommand{\QtopZtwoStwo}{--}
\newcommand{\QdieDownZtwoStwo}{--}
\newcommand{\NjointZtwoStwo}{528}
\newcommand{\PrtZtwoStwo}{97.5}
\newcommand{\UnderZtwoStwo}{37}
\newcommand{\RiseZtwoStwo}{56.1}
\newcommand{\PlossZthree}{120.3}
\newcommand{\EffZthree}{89.3}
\newcommand{\PconvZthree}{105.7}
\newcommand{\ProuteZthree}{14.63}
\newcommand{\PhvZthree}{--}
\newcommand{\PtsvZthree}{0.52}
\newcommand{\PattZthree}{2.29}
\newcommand{\PintraZthree}{1.48}
\newcommand{\PboardZthree}{10.29}
\newcommand{\PpadsZthree}{0.043}
\newcommand{\TmaxZthree}{82.4}
\newcommand{\TdieZthree}{78.5}
\newcommand{\TconvZthree}{45.9}
\newcommand{\TjointZthree}{78.7}
\newcommand{\WarpZthree}{4.7}
\newcommand{\PkgWarpZthree}{13.9}
\newcommand{\DuZthree}{1.41}
\newcommand{\EpsZthree}{1.98}
\newcommand{\EpsUZthree}{12.1}
\newcommand{\EtItersZthree}{4}
\newcommand{\TjointHotZthree}{79.2}
\newcommand{\TlsZthree}{50.3}
\newcommand{\TlsRiseZthree}{25}
\newcommand{\ThsZthree}{31.3}
\newcommand{\TindZthree}{56.3}
\newcommand{\QplateZthree}{173}
\newcommand{\QtopZthree}{947}
\newcommand{\QdieDownZthree}{53}
\newcommand{\NjointZthree}{1369}
\newcommand{\PrtZthree}{104.1}
\newcommand{\UnderZthree}{13}
\newcommand{\RiseZthree}{16.2}
\newcommand{\PlossZthreeSone}{121.8}
\newcommand{\EffZthreeSone}{89.1}
\newcommand{\PconvZthreeSone}{105.8}
\newcommand{\ProuteZthreeSone}{16.04}
\newcommand{\PhvZthreeSone}{--}
\newcommand{\PtsvZthreeSone}{0.52}
\newcommand{\PattZthreeSone}{2.29}
\newcommand{\PintraZthreeSone}{2.89}
\newcommand{\PboardZthreeSone}{10.29}
\newcommand{\PpadsZthreeSone}{0.043}
\newcommand{\TmaxZthreeSone}{82.5}
\newcommand{\TdieZthreeSone}{78.5}
\newcommand{\TconvZthreeSone}{46.0}
\newcommand{\TjointZthreeSone}{78.8}
\newcommand{\WarpZthreeSone}{4.8}
\newcommand{\PkgWarpZthreeSone}{14.0}
\newcommand{\DuZthreeSone}{1.41}
\newcommand{\EpsZthreeSone}{1.99}
\newcommand{\EpsUZthreeSone}{12.0}
\newcommand{\EtItersZthreeSone}{4}
\newcommand{\TjointHotZthreeSone}{79.2}
\newcommand{\TlsZthreeSone}{50.4}
\newcommand{\TlsRiseZthreeSone}{25}
\newcommand{\ThsZthreeSone}{31.3}
\newcommand{\TindZthreeSone}{56.4}
\newcommand{\QplateZthreeSone}{174}
\newcommand{\QtopZthreeSone}{948}
\newcommand{\QdieDownZthreeSone}{52}
\newcommand{\NjointZthreeSone}{1369}
\newcommand{\PrtZthreeSone}{105.4}
\newcommand{\UnderZthreeSone}{13}
\newcommand{\RiseZthreeSone}{16.4}
\newcommand{\PlossZthreeStwo}{125.0}
\newcommand{\EffZthreeStwo}{88.9}
\newcommand{\PconvZthreeStwo}{105.9}
\newcommand{\ProuteZthreeStwo}{19.05}
\newcommand{\PhvZthreeStwo}{--}
\newcommand{\PtsvZthreeStwo}{0.52}
\newcommand{\PattZthreeStwo}{2.30}
\newcommand{\PintraZthreeStwo}{5.89}
\newcommand{\PboardZthreeStwo}{10.30}
\newcommand{\PpadsZthreeStwo}{0.043}
\newcommand{\TmaxZthreeStwo}{82.5}
\newcommand{\TdieZthreeStwo}{78.6}
\newcommand{\TconvZthreeStwo}{46.3}
\newcommand{\TjointZthreeStwo}{78.9}
\newcommand{\WarpZthreeStwo}{4.9}
\newcommand{\PkgWarpZthreeStwo}{14.2}
\newcommand{\DuZthreeStwo}{1.42}
\newcommand{\EpsZthreeStwo}{2.02}
\newcommand{\EpsUZthreeStwo}{11.8}
\newcommand{\EtItersZthreeStwo}{4}
\newcommand{\TjointHotZthreeStwo}{79.3}
\newcommand{\TlsZthreeStwo}{50.7}
\newcommand{\TlsRiseZthreeStwo}{26}
\newcommand{\ThsZthreeStwo}{31.4}
\newcommand{\TindZthreeStwo}{56.7}
\newcommand{\QplateZthreeStwo}{176}
\newcommand{\QtopZthreeStwo}{948}
\newcommand{\QdieDownZthreeStwo}{51}
\newcommand{\NjointZthreeStwo}{1369}
\newcommand{\PrtZthreeStwo}{108.2}
\newcommand{\UnderZthreeStwo}{13}
\newcommand{\RiseZthreeStwo}{16.8}
\newcommand{\DeltaPZtwoSzero}{-3.13}
\newcommand{\DeltaPrtZtwoSzero}{0.59}
\newcommand{\DeltaPrtAbsZtwoSzero}{0.6}
\newcommand{\DeltaPZthreeSzero}{23.99}
\newcommand{\DeltaPrtZthreeSzero}{-10.02}
\newcommand{\DeltaPrtAbsZthreeSzero}{10.0}
\newcommand{\DeltaPZtwoSone}{30.51}
\newcommand{\DeltaPrtZtwoSone}{19.46}
\newcommand{\DeltaPrtAbsZtwoSone}{19.5}
\newcommand{\DeltaPZthreeSone}{58.13}
\newcommand{\DeltaPrtZthreeSone}{8.84}
\newcommand{\DeltaPrtAbsZthreeSone}{8.8}
\newcommand{\DeltaPZtwoStwo}{142.88}
\newcommand{\DeltaPrtZtwoStwo}{64.93}
\newcommand{\DeltaPrtAbsZtwoStwo}{64.9}
\newcommand{\DeltaPZthreeStwo}{171.59}
\newcommand{\DeltaPrtZthreeStwo}{54.32}
\newcommand{\DeltaPrtAbsZthreeStwo}{54.3}
\newcommand{\LossGapZtwoSzero}{3.1}
\newcommand{\LossLeadZtwoSzero}{Z1}
\newcommand{\dPconvZtwoSzero}{4.08}
\newcommand{\dPconvAbsZtwoSzero}{4.1}
\newcommand{\dProuteZtwoSzero}{-0.95}
\newcommand{\dProuteAbsZtwoSzero}{0.95}
\newcommand{\dTlsZtwoSzero}{7.0}
\newcommand{\dTlsAbsZtwoSzero}{7}
\newcommand{\dTmaxZtwoSzero}{8.7}
\newcommand{\dTmaxAbsZtwoSzero}{9}
\newcommand{\dTjointZtwoSzero}{-17.7}
\newcommand{\dTjointAbsZtwoSzero}{18}
\newcommand{\LossGapZthreeSzero}{24.0}
\newcommand{\LossLeadZthreeSzero}{Z3}
\newcommand{\dPconvZthreeSzero}{-34.99}
\newcommand{\dPconvAbsZthreeSzero}{35.0}
\newcommand{\dProuteZthreeSzero}{11.00}
\newcommand{\dProuteAbsZthreeSzero}{11.00}
\newcommand{\dTlsZthreeSzero}{-59.4}
\newcommand{\dTlsAbsZthreeSzero}{59}
\newcommand{\dTmaxZthreeSzero}{-33.2}
\newcommand{\dTmaxAbsZthreeSzero}{33}
\newcommand{\dTjointZthreeSzero}{-36.2}
\newcommand{\dTjointAbsZthreeSzero}{36}
\newcommand{\LossGapZtwoSone}{30.5}
\newcommand{\LossLeadZtwoSone}{Z2}
\newcommand{\dPconvZtwoSone}{0.20}
\newcommand{\dPconvAbsZtwoSone}{0.2}
\newcommand{\dProuteZtwoSone}{-30.71}
\newcommand{\dProuteAbsZtwoSone}{30.71}
\newcommand{\dTlsZtwoSone}{0.5}
\newcommand{\dTlsAbsZtwoSone}{0}
\newcommand{\dTmaxZtwoSone}{-15.7}
\newcommand{\dTmaxAbsZtwoSone}{16}
\newcommand{\dTjointZtwoSone}{-42.9}
\newcommand{\dTjointAbsZtwoSone}{43}
\newcommand{\LossGapZthreeSone}{58.1}
\newcommand{\LossLeadZthreeSone}{Z3}
\newcommand{\dPconvZthreeSone}{-39.04}
\newcommand{\dPconvAbsZthreeSone}{39.0}
\newcommand{\dProuteZthreeSone}{-19.09}
\newcommand{\dProuteAbsZthreeSone}{19.09}
\newcommand{\dTlsZthreeSone}{-66.3}
\newcommand{\dTlsAbsZthreeSone}{66}
\newcommand{\dTmaxZthreeSone}{-58.1}
\newcommand{\dTmaxAbsZthreeSone}{58}
\newcommand{\dTjointZthreeSone}{-61.5}
\newcommand{\dTjointAbsZthreeSone}{62}
\newcommand{\LossGapZtwoStwo}{142.9}
\newcommand{\LossLeadZtwoStwo}{Z2}
\newcommand{\dPconvZtwoStwo}{-12.83}
\newcommand{\dPconvAbsZtwoStwo}{12.8}
\newcommand{\dProuteZtwoStwo}{-130.05}
\newcommand{\dProuteAbsZtwoStwo}{130.05}
\newcommand{\dTlsZtwoStwo}{-21.6}
\newcommand{\dTlsAbsZtwoStwo}{22}
\newcommand{\dTmaxZtwoStwo}{-99.2}
\newcommand{\dTmaxAbsZtwoStwo}{99}
\newcommand{\dTjointZtwoStwo}{-126.5}
\newcommand{\dTjointAbsZtwoStwo}{127}
\newcommand{\LossGapZthreeStwo}{171.6}
\newcommand{\LossLeadZthreeStwo}{Z3}
\newcommand{\dPconvZthreeStwo}{-52.42}
\newcommand{\dPconvAbsZthreeStwo}{52.4}
\newcommand{\dProuteZthreeStwo}{-119.17}
\newcommand{\dProuteAbsZthreeStwo}{119.17}
\newcommand{\dTlsZthreeStwo}{-89.0}
\newcommand{\dTlsAbsZthreeStwo}{89}
\newcommand{\dTmaxZthreeStwo}{-142.6}
\newcommand{\dTmaxAbsZthreeStwo}{143}
\newcommand{\dTjointZthreeStwo}{-145.4}
\newcommand{\dTjointAbsZthreeStwo}{145}
\newcommand{\RiseSoneZone}{35.6}
\newcommand{\RiseSoneZtwo}{2.0}
\newcommand{\RiseSoneZthree}{1.5}
\newcommand{\ConvToRouteZtwo}{4.3}
\newcommand{\ConvToRouteZthree}{3.2}
\newcommand{\NvariantRuns}{57}
\newcommand{\WarpRatioZtwo}{2.69}
\newcommand{\WarpFoldZtwo}{0.4}
\newcommand{\DuRatioZtwo}{1.40}
\newcommand{\WarpRatioZthree}{0.19}
\newcommand{\WarpFoldZthree}{5.4}
\newcommand{\DuRatioZthree}{0.64}
\newcommand{\LandTie}{28}
\newcommand{\LandPT}{667}
\newcommand{\LandPmax}{67}
\newcommand{\LandTmin}{10}
\newcommand{\LandPmin}{11}
\newcommand{\LandTmax}{60}
\newcommand{\LandDVmax}{163}
\newcommand{\LandDVmin}{27}
\newcommand{\LandBallsInside}{245}
\newcommand{\LandBallsRail}{685}
\newcommand{\PitchMaxZthree}{0.58}
\newcommand{\SensTierRefZonePloss}{152.1}
\newcommand{\SensDTierRefZonePloss}{7.8}
\newcommand{\SensDAbsTierRefZonePloss}{7.8}
\newcommand{\SensTierRefZoneTmax}{128.3}
\newcommand{\SensDTierRefZoneTmax}{12.7}
\newcommand{\SensDAbsTierRefZoneTmax}{12.7}
\newcommand{\SensTierRefZoneTls}{123.3}
\newcommand{\SensDTierRefZoneTls}{13.6}
\newcommand{\SensDAbsTierRefZoneTls}{13.6}
\newcommand{\SensTierRefZoneWarp}{38.2}
\newcommand{\SensDTierRefZoneWarp}{12.7}
\newcommand{\SensDAbsTierRefZoneWarp}{12.7}
\newcommand{\SensTierRefZoneDu}{3.81}
\newcommand{\SensDTierRefZoneDu}{1.60}
\newcommand{\SensDAbsTierRefZoneDu}{1.60}
\newcommand{\SensTierRefZoneTjoint}{128.0}
\newcommand{\SensDTierRefZoneTjoint}{13.0}
\newcommand{\SensDAbsTierRefZoneTjoint}{13.0}
\newcommand{\SensZtwoCooledZtwoPloss}{111.6}
\newcommand{\SensDZtwoCooledZtwoPloss}{-35.9}
\newcommand{\SensDAbsZtwoCooledZtwoPloss}{35.9}
\newcommand{\SensZtwoCooledZtwoTmax}{80.1}
\newcommand{\SensDZtwoCooledZtwoTmax}{-44.3}
\newcommand{\SensDAbsZtwoCooledZtwoTmax}{44.3}
\newcommand{\SensZtwoCooledZtwoTls}{56.5}
\newcommand{\SensDZtwoCooledZtwoTls}{-60.2}
\newcommand{\SensDAbsZtwoCooledZtwoTls}{60.2}
\newcommand{\SensZtwoCooledZtwoWarp}{12.9}
\newcommand{\SensDZtwoCooledZtwoWarp}{-55.7}
\newcommand{\SensDAbsZtwoCooledZtwoWarp}{55.7}
\newcommand{\SensZtwoCooledZtwoDu}{2.00}
\newcommand{\SensDZtwoCooledZtwoDu}{-1.09}
\newcommand{\SensDAbsZtwoCooledZtwoDu}{1.09}
\newcommand{\SensZtwoCooledZtwoTjoint}{76.5}
\newcommand{\SensDZtwoCooledZtwoTjoint}{-20.7}
\newcommand{\SensDAbsZtwoCooledZtwoTjoint}{20.7}
\newcommand{\SensZthreeTopZthreePloss}{199.6}
\newcommand{\SensDZthreeTopZthreePloss}{79.2}
\newcommand{\SensDAbsZthreeTopZthreePloss}{79.2}
\newcommand{\SensZthreeTopZthreeTmax}{185.8}
\newcommand{\SensDZthreeTopZthreeTmax}{103.3}
\newcommand{\SensDAbsZthreeTopZthreeTmax}{103.3}
\newcommand{\SensZthreeTopZthreeTls}{179.7}
\newcommand{\SensDZthreeTopZthreeTls}{129.4}
\newcommand{\SensDAbsZthreeTopZthreeTls}{129.4}
\newcommand{\SensZthreeTopZthreeWarp}{59.5}
\newcommand{\SensDZthreeTopZthreeWarp}{54.7}
\newcommand{\SensDAbsZthreeTopZthreeWarp}{54.7}
\newcommand{\SensZthreeTopZthreeDu}{4.47}
\newcommand{\SensDZthreeTopZthreeDu}{3.07}
\newcommand{\SensDAbsZthreeTopZthreeDu}{3.07}
\newcommand{\SensZthreeTopZthreeTjoint}{109.9}
\newcommand{\SensDZthreeTopZthreeTjoint}{31.1}
\newcommand{\SensDAbsZthreeTopZthreeTjoint}{31.1}
\newcommand{\SensHbotLoZonePloss}{143.7}
\newcommand{\SensDHbotLoZonePloss}{-0.6}
\newcommand{\SensDAbsHbotLoZonePloss}{0.6}
\newcommand{\SensHbotLoZoneTmax}{114.1}
\newcommand{\SensDHbotLoZoneTmax}{-1.6}
\newcommand{\SensDAbsHbotLoZoneTmax}{1.6}
\newcommand{\SensHbotLoZoneTls}{108.8}
\newcommand{\SensDHbotLoZoneTls}{-0.9}
\newcommand{\SensDAbsHbotLoZoneTls}{0.9}
\newcommand{\SensHbotLoZoneWarp}{17.6}
\newcommand{\SensDHbotLoZoneWarp}{-8.0}
\newcommand{\SensDAbsHbotLoZoneWarp}{8.0}
\newcommand{\SensHbotLoZoneDu}{1.66}
\newcommand{\SensDHbotLoZoneDu}{-0.55}
\newcommand{\SensDAbsHbotLoZoneDu}{0.55}
\newcommand{\SensHbotLoZoneTjoint}{111.3}
\newcommand{\SensDHbotLoZoneTjoint}{-3.7}
\newcommand{\SensDAbsHbotLoZoneTjoint}{3.7}
\newcommand{\SensHbotLoZtwoPloss}{147.4}
\newcommand{\SensDHbotLoZtwoPloss}{-0.1}
\newcommand{\SensDAbsHbotLoZtwoPloss}{0.1}
\newcommand{\SensHbotLoZtwoTmax}{124.0}
\newcommand{\SensDHbotLoZtwoTmax}{-0.3}
\newcommand{\SensDAbsHbotLoZtwoTmax}{0.3}
\newcommand{\SensHbotLoZtwoTls}{116.6}
\newcommand{\SensDHbotLoZtwoTls}{-0.1}
\newcommand{\SensDAbsHbotLoZtwoTls}{0.1}
\newcommand{\SensHbotLoZtwoWarp}{65.3}
\newcommand{\SensDHbotLoZtwoWarp}{-3.4}
\newcommand{\SensDAbsHbotLoZtwoWarp}{3.4}
\newcommand{\SensHbotLoZtwoDu}{2.57}
\newcommand{\SensDHbotLoZtwoDu}{-0.52}
\newcommand{\SensDAbsHbotLoZtwoDu}{0.52}
\newcommand{\SensHbotLoZtwoTjoint}{95.4}
\newcommand{\SensDHbotLoZtwoTjoint}{-1.8}
\newcommand{\SensDAbsHbotLoZtwoTjoint}{1.8}
\newcommand{\SensHbotHiZonePloss}{143.0}
\newcommand{\SensDHbotHiZonePloss}{-1.3}
\newcommand{\SensDAbsHbotHiZonePloss}{1.3}
\newcommand{\SensHbotHiZoneTmax}{112.2}
\newcommand{\SensDHbotHiZoneTmax}{-3.4}
\newcommand{\SensDAbsHbotHiZoneTmax}{3.4}
\newcommand{\SensHbotHiZoneTls}{107.6}
\newcommand{\SensDHbotHiZoneTls}{-2.1}
\newcommand{\SensDAbsHbotHiZoneTls}{2.1}
\newcommand{\SensHbotHiZoneWarp}{7.7}
\newcommand{\SensDHbotHiZoneWarp}{-17.8}
\newcommand{\SensDAbsHbotHiZoneWarp}{17.8}
\newcommand{\SensHbotHiZoneDu}{1.01}
\newcommand{\SensDHbotHiZoneDu}{-1.20}
\newcommand{\SensDAbsHbotHiZoneDu}{1.20}
\newcommand{\SensHbotHiZoneTjoint}{102.7}
\newcommand{\SensDHbotHiZoneTjoint}{-12.3}
\newcommand{\SensDAbsHbotHiZoneTjoint}{12.3}
\newcommand{\SensHbotHiZtwoPloss}{147.3}
\newcommand{\SensDHbotHiZtwoPloss}{-0.2}
\newcommand{\SensDAbsHbotHiZtwoPloss}{0.2}
\newcommand{\SensHbotHiZtwoTmax}{123.6}
\newcommand{\SensDHbotHiZtwoTmax}{-0.7}
\newcommand{\SensDAbsHbotHiZtwoTmax}{0.7}
\newcommand{\SensHbotHiZtwoTls}{116.5}
\newcommand{\SensDHbotHiZtwoTls}{-0.3}
\newcommand{\SensDAbsHbotHiZtwoTls}{0.3}
\newcommand{\SensHbotHiZtwoWarp}{61.2}
\newcommand{\SensDHbotHiZtwoWarp}{-7.4}
\newcommand{\SensDAbsHbotHiZtwoWarp}{7.4}
\newcommand{\SensHbotHiZtwoDu}{1.93}
\newcommand{\SensDHbotHiZtwoDu}{-1.15}
\newcommand{\SensDAbsHbotHiZtwoDu}{1.15}
\newcommand{\SensHbotHiZtwoTjoint}{93.3}
\newcommand{\SensDHbotHiZtwoTjoint}{-3.9}
\newcommand{\SensDAbsHbotHiZtwoTjoint}{3.9}
\newcommand{\SensTsvMinZonePloss}{144.3}
\newcommand{\SensDTsvMinZonePloss}{0.0}
\newcommand{\SensDAbsTsvMinZonePloss}{0.0}
\newcommand{\SensTsvMinZoneTmax}{115.7}
\newcommand{\SensDTsvMinZoneTmax}{0.0}
\newcommand{\SensDAbsTsvMinZoneTmax}{0.0}
\newcommand{\SensTsvMinZoneTls}{109.7}
\newcommand{\SensDTsvMinZoneTls}{0.0}
\newcommand{\SensDAbsTsvMinZoneTls}{0.0}
\newcommand{\SensTsvMinZoneWarp}{25.4}
\newcommand{\SensDTsvMinZoneWarp}{-0.2}
\newcommand{\SensDAbsTsvMinZoneWarp}{0.2}
\newcommand{\SensTsvMinZoneDu}{2.22}
\newcommand{\SensDTsvMinZoneDu}{0.01}
\newcommand{\SensDAbsTsvMinZoneDu}{0.01}
\newcommand{\SensTsvMinZoneTjoint}{115.0}
\newcommand{\SensDTsvMinZoneTjoint}{0.0}
\newcommand{\SensDAbsTsvMinZoneTjoint}{0.0}
\newcommand{\SensTsvMinZtwoPloss}{160.9}
\newcommand{\SensDTsvMinZtwoPloss}{13.5}
\newcommand{\SensDAbsTsvMinZtwoPloss}{13.5}
\newcommand{\SensTsvMinZtwoTmax}{126.0}
\newcommand{\SensDTsvMinZtwoTmax}{1.7}
\newcommand{\SensDAbsTsvMinZtwoTmax}{1.7}
\newcommand{\SensTsvMinZtwoTls}{118.4}
\newcommand{\SensDTsvMinZtwoTls}{1.7}
\newcommand{\SensDAbsTsvMinZtwoTls}{1.7}
\newcommand{\SensTsvMinZtwoWarp}{71.7}
\newcommand{\SensDTsvMinZtwoWarp}{3.0}
\newcommand{\SensDAbsTsvMinZtwoWarp}{3.0}
\newcommand{\SensTsvMinZtwoDu}{3.19}
\newcommand{\SensDTsvMinZtwoDu}{0.11}
\newcommand{\SensDAbsTsvMinZtwoDu}{0.11}
\newcommand{\SensTsvMinZtwoTjoint}{98.3}
\newcommand{\SensDTsvMinZtwoTjoint}{1.1}
\newcommand{\SensDAbsTsvMinZtwoTjoint}{1.1}
\newcommand{\SensTsvMidZonePloss}{144.3}
\newcommand{\SensDTsvMidZonePloss}{0.0}
\newcommand{\SensDAbsTsvMidZonePloss}{0.0}
\newcommand{\SensTsvMidZoneTmax}{115.7}
\newcommand{\SensDTsvMidZoneTmax}{0.0}
\newcommand{\SensDAbsTsvMidZoneTmax}{0.0}
\newcommand{\SensTsvMidZoneTls}{109.7}
\newcommand{\SensDTsvMidZoneTls}{0.0}
\newcommand{\SensDAbsTsvMidZoneTls}{0.0}
\newcommand{\SensTsvMidZoneWarp}{25.5}
\newcommand{\SensDTsvMidZoneWarp}{-0.1}
\newcommand{\SensDAbsTsvMidZoneWarp}{0.1}
\newcommand{\SensTsvMidZoneDu}{2.21}
\newcommand{\SensDTsvMidZoneDu}{0.01}
\newcommand{\SensDAbsTsvMidZoneDu}{0.01}
\newcommand{\SensTsvMidZoneTjoint}{115.0}
\newcommand{\SensDTsvMidZoneTjoint}{0.0}
\newcommand{\SensDAbsTsvMidZoneTjoint}{0.0}
\newcommand{\SensTsvMidZtwoPloss}{148.1}
\newcommand{\SensDTsvMidZtwoPloss}{0.7}
\newcommand{\SensDAbsTsvMidZtwoPloss}{0.7}
\newcommand{\SensTsvMidZtwoTmax}{124.5}
\newcommand{\SensDTsvMidZtwoTmax}{0.2}
\newcommand{\SensDAbsTsvMidZtwoTmax}{0.2}
\newcommand{\SensTsvMidZtwoTls}{116.9}
\newcommand{\SensDTsvMidZtwoTls}{0.2}
\newcommand{\SensDAbsTsvMidZtwoTls}{0.2}
\newcommand{\SensTsvMidZtwoWarp}{69.7}
\newcommand{\SensDTsvMidZtwoWarp}{1.0}
\newcommand{\SensDAbsTsvMidZtwoWarp}{1.0}
\newcommand{\SensTsvMidZtwoDu}{3.12}
\newcommand{\SensDTsvMidZtwoDu}{0.03}
\newcommand{\SensDAbsTsvMidZtwoDu}{0.03}
\newcommand{\SensTsvMidZtwoTjoint}{97.2}
\newcommand{\SensDTsvMidZtwoTjoint}{0.0}
\newcommand{\SensDAbsTsvMidZtwoTjoint}{0.0}
\newcommand{\SensAttSparseZtwoPloss}{148.1}
\newcommand{\SensDAttSparseZtwoPloss}{0.7}
\newcommand{\SensDAbsAttSparseZtwoPloss}{0.7}
\newcommand{\SensAttSparseZtwoTmax}{125.1}
\newcommand{\SensDAttSparseZtwoTmax}{0.8}
\newcommand{\SensDAbsAttSparseZtwoTmax}{0.8}
\newcommand{\SensAttSparseZtwoTls}{117.5}
\newcommand{\SensDAttSparseZtwoTls}{0.8}
\newcommand{\SensDAbsAttSparseZtwoTls}{0.8}
\newcommand{\SensAttSparseZtwoWarp}{69.2}
\newcommand{\SensDAttSparseZtwoWarp}{0.6}
\newcommand{\SensDAbsAttSparseZtwoWarp}{0.6}
\newcommand{\SensAttSparseZtwoDu}{3.09}
\newcommand{\SensDAttSparseZtwoDu}{0.00}
\newcommand{\SensDAbsAttSparseZtwoDu}{0.00}
\newcommand{\SensAttSparseZtwoTjoint}{97.3}
\newcommand{\SensDAttSparseZtwoTjoint}{0.1}
\newcommand{\SensDAbsAttSparseZtwoTjoint}{0.1}
\newcommand{\SensVoidZonePloss}{146.5}
\newcommand{\SensDVoidZonePloss}{2.2}
\newcommand{\SensDAbsVoidZonePloss}{2.2}
\newcommand{\SensVoidZoneTmax}{119.9}
\newcommand{\SensDVoidZoneTmax}{4.3}
\newcommand{\SensDAbsVoidZoneTmax}{4.3}
\newcommand{\SensVoidZoneTls}{113.4}
\newcommand{\SensDVoidZoneTls}{3.7}
\newcommand{\SensDAbsVoidZoneTls}{3.7}
\newcommand{\SensVoidZoneWarp}{34.6}
\newcommand{\SensDVoidZoneWarp}{9.1}
\newcommand{\SensDAbsVoidZoneWarp}{9.1}
\newcommand{\SensVoidZoneDu}{3.48}
\newcommand{\SensDVoidZoneDu}{1.27}
\newcommand{\SensDAbsVoidZoneDu}{1.27}
\newcommand{\SensVoidZoneTjoint}{119.5}
\newcommand{\SensDVoidZoneTjoint}{4.6}
\newcommand{\SensDAbsVoidZoneTjoint}{4.6}
\newcommand{\SensVoidZtwoPloss}{149.3}
\newcommand{\SensDVoidZtwoPloss}{1.8}
\newcommand{\SensDAbsVoidZtwoPloss}{1.8}
\newcommand{\SensVoidZtwoTmax}{127.2}
\newcommand{\SensDVoidZtwoTmax}{2.9}
\newcommand{\SensDAbsVoidZtwoTmax}{2.9}
\newcommand{\SensVoidZtwoTls}{119.9}
\newcommand{\SensDVoidZtwoTls}{3.2}
\newcommand{\SensDAbsVoidZtwoTls}{3.2}
\newcommand{\SensVoidZtwoWarp}{39.1}
\newcommand{\SensDVoidZtwoWarp}{-29.5}
\newcommand{\SensDAbsVoidZtwoWarp}{29.5}
\newcommand{\SensVoidZtwoDu}{3.03}
\newcommand{\SensDVoidZtwoDu}{-0.06}
\newcommand{\SensDAbsVoidZtwoDu}{0.06}
\newcommand{\SensVoidZtwoTjoint}{97.5}
\newcommand{\SensDVoidZtwoTjoint}{0.3}
\newcommand{\SensDAbsVoidZtwoTjoint}{0.3}
\newcommand{\SensIcepakZonePloss}{142.4}
\newcommand{\SensDIcepakZonePloss}{-1.9}
\newcommand{\SensDAbsIcepakZonePloss}{1.9}
\newcommand{\SensIcepakZoneTmax}{112.5}
\newcommand{\SensDIcepakZoneTmax}{-3.2}
\newcommand{\SensDAbsIcepakZoneTmax}{3.2}
\newcommand{\SensIcepakZoneTls}{106.6}
\newcommand{\SensDIcepakZoneTls}{-3.2}
\newcommand{\SensDAbsIcepakZoneTls}{3.2}
\newcommand{\SensIcepakZoneWarp}{24.7}
\newcommand{\SensDIcepakZoneWarp}{-0.8}
\newcommand{\SensDAbsIcepakZoneWarp}{0.8}
\newcommand{\SensIcepakZoneDu}{2.13}
\newcommand{\SensDIcepakZoneDu}{-0.08}
\newcommand{\SensDAbsIcepakZoneDu}{0.08}
\newcommand{\SensIcepakZoneTjoint}{111.8}
\newcommand{\SensDIcepakZoneTjoint}{-3.2}
\newcommand{\SensDAbsIcepakZoneTjoint}{3.2}
\newcommand{\SensIcepakZtwoPloss}{145.5}
\newcommand{\SensDIcepakZtwoPloss}{-1.9}
\newcommand{\SensDAbsIcepakZtwoPloss}{1.9}
\newcommand{\SensIcepakZtwoTmax}{121.1}
\newcommand{\SensDIcepakZtwoTmax}{-3.3}
\newcommand{\SensDAbsIcepakZtwoTmax}{3.3}
\newcommand{\SensIcepakZtwoTls}{113.5}
\newcommand{\SensDIcepakZtwoTls}{-3.3}
\newcommand{\SensDAbsIcepakZtwoTls}{3.3}
\newcommand{\SensIcepakZtwoWarp}{66.3}
\newcommand{\SensDIcepakZtwoWarp}{-2.4}
\newcommand{\SensDAbsIcepakZtwoWarp}{2.4}
\newcommand{\SensIcepakZtwoDu}{2.96}
\newcommand{\SensDIcepakZtwoDu}{-0.13}
\newcommand{\SensDAbsIcepakZtwoDu}{0.13}
\newcommand{\SensIcepakZtwoTjoint}{94.3}
\newcommand{\SensDIcepakZtwoTjoint}{-2.9}
\newcommand{\SensDAbsIcepakZtwoTjoint}{2.9}
\newcommand{\SensGanHalfZonePloss}{144.8}
\newcommand{\SensDGanHalfZonePloss}{0.5}
\newcommand{\SensDAbsGanHalfZonePloss}{0.5}
\newcommand{\SensGanHalfZoneTmax}{117.2}
\newcommand{\SensDGanHalfZoneTmax}{1.6}
\newcommand{\SensDAbsGanHalfZoneTmax}{1.6}
\newcommand{\SensGanHalfZoneTls}{110.6}
\newcommand{\SensDGanHalfZoneTls}{0.8}
\newcommand{\SensDAbsGanHalfZoneTls}{0.8}
\newcommand{\SensGanHalfZoneWarp}{26.2}
\newcommand{\SensDGanHalfZoneWarp}{0.6}
\newcommand{\SensDAbsGanHalfZoneWarp}{0.6}
\newcommand{\SensGanHalfZoneDu}{2.26}
\newcommand{\SensDGanHalfZoneDu}{0.05}
\newcommand{\SensDAbsGanHalfZoneDu}{0.05}
\newcommand{\SensGanHalfZoneTjoint}{116.5}
\newcommand{\SensDGanHalfZoneTjoint}{1.5}
\newcommand{\SensDAbsGanHalfZoneTjoint}{1.5}
\newcommand{\SensGanHalfZtwoPloss}{148.0}
\newcommand{\SensDGanHalfZtwoPloss}{0.5}
\newcommand{\SensDAbsGanHalfZtwoPloss}{0.5}
\newcommand{\SensGanHalfZtwoTmax}{126.1}
\newcommand{\SensDGanHalfZtwoTmax}{1.8}
\newcommand{\SensDAbsGanHalfZtwoTmax}{1.8}
\newcommand{\SensGanHalfZtwoTls}{117.6}
\newcommand{\SensDGanHalfZtwoTls}{0.9}
\newcommand{\SensDAbsGanHalfZtwoTls}{0.9}
\newcommand{\SensGanHalfZtwoWarp}{69.6}
\newcommand{\SensDGanHalfZtwoWarp}{0.9}
\newcommand{\SensDAbsGanHalfZtwoWarp}{0.9}
\newcommand{\SensGanHalfZtwoDu}{3.10}
\newcommand{\SensDGanHalfZtwoDu}{0.01}
\newcommand{\SensDAbsGanHalfZtwoDu}{0.01}
\newcommand{\SensGanHalfZtwoTjoint}{97.3}
\newcommand{\SensDGanHalfZtwoTjoint}{0.1}
\newcommand{\SensDAbsGanHalfZtwoTjoint}{0.1}
\newcommand{\SensBoardHoneZonePloss}{142.4}
\newcommand{\SensDBoardHoneZonePloss}{-1.9}
\newcommand{\SensDAbsBoardHoneZonePloss}{1.9}
\newcommand{\SensBoardHoneZoneTmax}{112.6}
\newcommand{\SensDBoardHoneZoneTmax}{-3.1}
\newcommand{\SensDAbsBoardHoneZoneTmax}{3.1}
\newcommand{\SensBoardHoneZoneTls}{106.6}
\newcommand{\SensDBoardHoneZoneTls}{-3.2}
\newcommand{\SensDAbsBoardHoneZoneTls}{3.2}
\newcommand{\SensBoardHoneZoneWarp}{24.3}
\newcommand{\SensDBoardHoneZoneWarp}{-1.3}
\newcommand{\SensDAbsBoardHoneZoneWarp}{1.3}
\newcommand{\SensBoardHoneZoneDu}{2.08}
\newcommand{\SensDBoardHoneZoneDu}{-0.13}
\newcommand{\SensDAbsBoardHoneZoneDu}{0.13}
\newcommand{\SensBoardHoneZoneTjoint}{107.5}
\newcommand{\SensDBoardHoneZoneTjoint}{-7.4}
\newcommand{\SensDAbsBoardHoneZoneTjoint}{7.4}
\newcommand{\SensBoardHoneZtwoPloss}{147.3}
\newcommand{\SensDBoardHoneZtwoPloss}{-0.2}
\newcommand{\SensDAbsBoardHoneZtwoPloss}{0.2}
\newcommand{\SensBoardHoneZtwoTmax}{124.0}
\newcommand{\SensDBoardHoneZtwoTmax}{-0.4}
\newcommand{\SensDAbsBoardHoneZtwoTmax}{0.4}
\newcommand{\SensBoardHoneZtwoTls}{116.4}
\newcommand{\SensDBoardHoneZtwoTls}{-0.3}
\newcommand{\SensDAbsBoardHoneZtwoTls}{0.3}
\newcommand{\SensBoardHoneZtwoWarp}{65.6}
\newcommand{\SensDBoardHoneZtwoWarp}{-3.0}
\newcommand{\SensDAbsBoardHoneZtwoWarp}{3.0}
\newcommand{\SensBoardHoneZtwoDu}{2.56}
\newcommand{\SensDBoardHoneZtwoDu}{-0.53}
\newcommand{\SensDAbsBoardHoneZtwoDu}{0.53}
\newcommand{\SensBoardHoneZtwoTjoint}{95.4}
\newcommand{\SensDBoardHoneZtwoTjoint}{-1.8}
\newcommand{\SensDAbsBoardHoneZtwoTjoint}{1.8}
\newcommand{\PlossAZone}{144.3}
\newcommand{\PconvAZone}{140.7}
\newcommand{\ProuteAZone}{3.6}
\newcommand{\TmaxAZone}{115.7}
\newcommand{\TdieAZone}{89.8}
\newcommand{\TlsAZone}{109.7}
\newcommand{\WarpAZone}{25.5}
\newcommand{\DuAZone}{2.21}
\newcommand{\TjointAZone}{115.0}
\newcommand{\TjointHotAZone}{115.4}
\newcommand{\EpsAZone}{5.80}
\newcommand{\PlossAZtwo}{147.4}
\newcommand{\PconvAZtwo}{144.8}
\newcommand{\ProuteAZtwo}{2.7}
\newcommand{\TmaxAZtwo}{124.3}
\newcommand{\TdieAZtwo}{90.0}
\newcommand{\TlsAZtwo}{116.7}
\newcommand{\WarpAZtwo}{68.7}
\newcommand{\DuAZtwo}{3.09}
\newcommand{\TjointAZtwo}{97.2}
\newcommand{\TjointHotAZtwo}{98.3}
\newcommand{\EpsAZtwo}{8.03}
\newcommand{\PlossAZthree}{199.6}
\newcommand{\PconvAZthree}{181.4}
\newcommand{\ProuteAZthree}{18.1}
\newcommand{\TmaxAZthree}{185.8}
\newcommand{\TdieAZthree}{93.0}
\newcommand{\TlsAZthree}{179.7}
\newcommand{\WarpAZthree}{59.5}
\newcommand{\DuAZthree}{4.47}
\newcommand{\TjointAZthree}{109.9}
\newcommand{\TjointHotAZthree}{118.6}
\newcommand{\QbotAZthree}{1}
\newcommand{\EpsAZthree}{13.47}
\newcommand{\PlossBZone}{142.6}
\newcommand{\PconvBZone}{139.1}
\newcommand{\ProuteBZone}{3.6}
\newcommand{\TmaxBZone}{111.3}
\newcommand{\TdieBZone}{89.0}
\newcommand{\TlsBZone}{107.0}
\newcommand{\WarpBZone}{5.9}
\newcommand{\DuBZone}{0.69}
\newcommand{\TjointBZone}{108.3}
\newcommand{\TjointHotBZone}{110.1}
\newcommand{\QbotBZone}{13}
\newcommand{\EpsBZone}{2.08}
\newcommand{\PlossBZtwo}{115.2}
\newcommand{\PconvBZtwo}{112.9}
\newcommand{\ProuteBZtwo}{2.3}
\newcommand{\TmaxBZtwo}{81.5}
\newcommand{\TdieBZtwo}{77.2}
\newcommand{\TlsBZtwo}{62.6}
\newcommand{\WarpBZtwo}{17.6}
\newcommand{\DuBZtwo}{2.10}
\newcommand{\TjointBZtwo}{78.7}
\newcommand{\TjointHotBZtwo}{79.1}
\newcommand{\QbotBZtwo}{190}
\newcommand{\EpsBZtwo}{4.15}
\newcommand{\PlossBZthree}{124.5}
\newcommand{\PconvBZthree}{109.7}
\newcommand{\ProuteBZthree}{14.8}
\newcommand{\TmaxBZthree}{83.3}
\newcommand{\TdieBZthree}{79.3}
\newcommand{\TlsBZthree}{57.0}
\newcommand{\WarpBZthree}{7.1}
\newcommand{\DuBZthree}{1.56}
\newcommand{\TjointBZthree}{80.4}
\newcommand{\TjointHotBZthree}{80.7}
\newcommand{\QbotBZthree}{164}
\newcommand{\EpsBZthree}{2.52}
\newcommand{\PlossIZtwo}{111.6}
\newcommand{\PconvIZtwo}{109.3}
\newcommand{\ProuteIZtwo}{2.3}
\newcommand{\TmaxIZtwo}{80.1}
\newcommand{\TdieIZtwo}{75.8}
\newcommand{\TlsIZtwo}{56.5}
\newcommand{\WarpIZtwo}{12.9}
\newcommand{\DuIZtwo}{2.00}
\newcommand{\TjointIZtwo}{76.5}
\newcommand{\TjointHotIZtwo}{76.9}
\newcommand{\QbotIZtwo}{212}
\newcommand{\EpsIZtwo}{3.72}
\newcommand{\PlossIZthree}{120.3}
\newcommand{\PconvIZthree}{105.7}
\newcommand{\ProuteIZthree}{14.6}
\newcommand{\TmaxIZthree}{82.4}
\newcommand{\TdieIZthree}{78.5}
\newcommand{\TlsIZthree}{50.3}
\newcommand{\WarpIZthree}{4.7}
\newcommand{\DuIZthree}{1.41}
\newcommand{\TjointIZthree}{78.7}
\newcommand{\TjointHotIZthree}{79.2}
\newcommand{\QbotIZthree}{173}
\newcommand{\EpsIZthree}{2.04}
\newcommand{\GapAZtwo}{3.1}
\newcommand{\LeadAZtwo}{Z1}
\newcommand{\dTmaxTierAZtwo}{9}
\newcommand{\GapAZthree}{55.2}
\newcommand{\LeadAZthree}{Z1}
\newcommand{\dTmaxTierAZthree}{70}
\newcommand{\GapBZtwo}{27.4}
\newcommand{\LeadBZtwo}{Z2}
\newcommand{\dTmaxTierBZtwo}{30}
\newcommand{\GapBZthree}{18.2}
\newcommand{\LeadBZthree}{Z3}
\newcommand{\dTmaxTierBZthree}{28}
\newcommand{\GainBZone}{1.7}
\newcommand{\CoolBZone}{4}
\newcommand{\GainBZtwo}{32.3}
\newcommand{\CoolBZtwo}{43}
\newcommand{\GainBZthree}{75.1}
\newcommand{\CoolBZthree}{102}
\newcommand{\PIsoZone}{94.1}
\newcommand{\ProuteIsoZone}{2.6}
\newcommand{\PconvIso}{91.5}
\newcommand{\PIsoZtwo}{93.5}
\newcommand{\ProuteIsoZtwo}{2.0}
\newcommand{\PIsoZthree}{104.1}
\newcommand{\ProuteIsoZthree}{12.6}
\newcommand{\GapIsoZtwo}{0.6}
\newcommand{\GapIsoZthree}{10.0}
\newcommand{\PIsoHotZone}{138.8}
\newcommand{\ProuteIsoHotZone}{3.5}
\newcommand{\PconvIsoHot}{135.3}
\newcommand{\PIsoHotZtwo}{137.9}
\newcommand{\ProuteIsoHotZtwo}{2.6}
\newcommand{\PIsoHotZthree}{151.6}
\newcommand{\ProuteIsoHotZthree}{16.3}
\newcommand{\GapIsoHotZtwo}{0.9}
\newcommand{\GapIsoHotZthree}{12.8}
\newcommand{\TboardAZone}{111}
\newcommand{\TboardBotAZone}{108}
\newcommand{\TsubAZone}{113}
\newcommand{\CteBoard}{14.7}
\newcommand{\CteSi}{2.6}
\newcommand{\TboardBZone}{67}
\newcommand{\TboardBotBZone}{28}
\newcommand{\TsubBZone}{107}
\newcommand{\TboardAZtwo}{93}
\newcommand{\TboardBotAZtwo}{91}
\newcommand{\TsubAZtwo}{95}
\newcommand{\TboardBZtwo}{75}
\newcommand{\TboardBotBZtwo}{74}
\newcommand{\TsubBZtwo}{76}
\newcommand{\TboardAZthree}{131}
\newcommand{\TboardBotAZthree}{149}
\newcommand{\TsubAZthree}{102}
\newcommand{\TboardBZthree}{72}
\newcommand{\TboardBotBZthree}{68}
\newcommand{\TsubBZthree}{79}
\newcommand{\MeshZthreePloss}{0.2}
\newcommand{\MeshZthreeTmax}{0.1}
\newcommand{\MeshZthreeWarp}{3.2}
\newcommand{\MeshZthreeDu}{1.1}
\newcommand{\RstarTierAZtwo}{none}
\newcommand{\RstarTierAZthree}{none}
\newcommand{\RstarTierBZtwo}{48.4}
\newcommand{\RstarTierBZthree}{32.1}
\newcommand{\RstarZtwoSzero}{none}
\newcommand{\AdvZtwoSzero}{-3.13}
\newcommand{\IrmsSzero}{23.3}
\newcommand{\PfeedPerMohmSzero}{0.54}
\newcommand{\RstarZtwoSone}{3.16}
\newcommand{\AdvZtwoSone}{30.51}
\newcommand{\IrmsSone}{93.5}
\newcommand{\PfeedPerMohmSone}{8.74}
\newcommand{\RstarZtwoStwo}{--}
\newcommand{\AdvZtwoStwo}{142.88}
\newcommand{\IrmsStwo}{162.4}
\newcommand{\PfeedPerMohmStwo}{26.37}
\newcommand{\RstarZthreeSzero}{42.22}
\newcommand{\AdvZthreeSzero}{23.99}
\newcommand{\RstarZthreeSone}{6.01}
\newcommand{\AdvZthreeSone}{58.13}
\newcommand{\RstarZthreeStwo}{--}
\newcommand{\AdvZthreeStwo}{171.59}
\newcommand{\MeshWorstProd}{1.1}
\newcommand{\LocalUrange}{0.8--12.1\,\%}
\newcommand{\EvalHours}{1.3--2.8~h}
\newcommand{\EnergyTol}{1e-06}
\newcommand{\EtItersMax}{7}
\newcommand{\NballMatched}{54}
\newcommand{\RloopZone}{2.26}
\newcommand{\RloopZtwo}{1.54}
\newcommand{\RstarOverLoopZtwo}{1.4}
\newcommand{\ErrTsvBulk}{1.1}
\newcommand{\ErrTsvKz}{0.71}
\newcommand{\ErrTsvKx}{0.25}
\newcommand{\ErrRdlKx}{3.7}
\newcommand{\ErrRdlKz}{12.5}
\newcommand{\ErrPkgRdlZone}{0.08}
\newcommand{\ErrThermAedtZone}{0.12}
\newcommand{\ErrThermAedtZtwo}{0.11}
\newcommand{\ErrThermMooseZone}{0.13}
\newcommand{\ErrThermMooseZtwo}{0.12}
\newcommand{\ErrMechClosed}{0.43}
\newcommand{\ErrMechAedt}{0.01}
\newcommand{\ErrCoolingIcepak}{16.5}
\newcommand{\ErrAttDenseKz}{0.49}
\newcommand{\ErrAttDenseKxy}{0.36}
\newcommand{\ErrHomogMax}{1.3}
\newcommand{\ErrThermMax}{0.1}

\newcommand{\pending}[1]{\textcolor{red}{[PENDING: #1]}}
\newenvironment{comment}{\par\medskip\noindent\textbf{Comment.}\itshape}{\par}
\newenvironment{reply}{\par\smallskip\noindent\textbf{Response.}}{\par\medskip}
\setlist{nosep,leftmargin=*}
\title{Response to the Reviewer\\[2pt]\large ``From Device to Solder Joint: An
Electrothermal--Mechanical Framework for Co-Designing Voltage-Regulator Integration in
Distributed Vertical Power Delivery''}
\author{Salma Abdelzaher, Mohamed Gharib, and Inna Partin-Vaisband}
\date{}
\begin{document}
\maketitle

We thank the reviewer for a careful and constructive report. The concerns about loss accounting,
component ratings, decoupling, the critical-joint selection and the uncertainty of close
comparisons were all well founded, and each is answered below with a new analysis rather than a
change of wording. Every number in this letter is generated from the same result files as the
revised manuscript. Section, equation, figure and table numbers refer to the revised manuscript;
``S'' tables are in the new supplementary material, which lists every material card, layer, mesh,
boundary condition, solver setting, runtime, variant and uncertainty.

\paragraph{Summary of the main changes.}
\begin{itemize}
\item \textbf{Z3 loss accounting (MC3).} The land-side redistribution of the 1-V current is now a
term of the coupled ledger, deposited in the substrate's land-side layer and iterated with the
temperature field; Z3 is re-solved for several land-side copper thicknesses, and its ranking, feed
crossover and voltage drop are re-evaluated. \pending{z3land numbers}
\item \textbf{Input decoupling (MC4).} The embedded design's local input capacitance is
implemented with real parts, verified by a circuit simulation including parasitics, and its effect
on the package network and the coupled solution is evaluated. \pending{decap}
\item \textbf{Feasibility filter (MC5).} A feasibility stage checks the GaN junction temperature,
capacitor, inductor, dielectric and encapsulant limits, the solder law's range and the
interconnect rules for every case; candidates that fail are removed before any ranking. Z3 under
the top plate alone fails six limits (junction bound \TjAZthree\,$^\circ$C against 150\,$^\circ$C).
\item \textbf{Omitted converter losses (MC6)} are bounded from datasheets at each case's
temperatures: \OmitLo--\OmitHi~W per integration; no ranking changes, and the count optimum is
unchanged.
\item \textbf{Count optimum through the full chain (MC7).} \pending{countopt}
\item \textbf{Solder joint (MC8, MC9).} \pending{solder: Pareto-set verification, 16 cycles, band,
alternative law, parameter table}
\item \textbf{Coupling (MC2).} Each coupling direction is classified (Section~II, new
Fig.~2/Algorithm); the one first-order mechanical-to-thermal path---the top interface following a
warped die---is implemented as an option and quantified.
\item \textbf{Cooling (MC10).} Complete cold-plate specification, entrance-length check, pressure
drop (\CoolDp~kPa) and pumping power (\CoolPhyd~W hydraulic per plate), heat balance per plate,
a non-uniform die power map, and the correction of a conservative double count of the caloric
resistance that this review led us to find.
\item \textbf{Verification and uncertainty (MC12, MC13).} A complete-package cross-check against
MOOSE on the identical mesh (agreement \FullPkgErrZone~K), refined MOOSE solutions that revised our
discretization estimate upward, and a per-output uncertainty budget with an explicit rule
(Table~V of the manuscript).
\item \textbf{Feeds and platform accounting (MC16).} Two physical feeds are evaluated against the
crossover resistance, and platform-inclusive losses including pumping power are reported.
\item \textbf{Positioning (MC1).} A comparison table against the closest work, a precise
statement of what is adopted, extended and new, and the removal of categorical claims.
\end{itemize}

\section*{Major comments}

\begin{comment}
MC1 -- The novelty relative to existing DVPD frameworks must be established more precisely;
comparison table; ablation.
\end{comment}
\begin{reply}
We agree, and have rewritten the related-work and contribution statements (Section~I). A new
Table~I compares the closest works---the DVPD design framework [1], the self-consistent
electrothermal model [18], and commercial multiphysics sign-off tools---with this work, capability
by capability; every cell for the other works is supported by a quoted passage (the detailed version
with pointers is Table~S1). The text now states explicitly:
(i) \emph{adopted from [1]}: the 1-kW specification, the embedded integration, the TSV and BGA
technology, the routing model used in the count screen, and the input-current model used as
scenario S2; \emph{adopted from our converter work [23], [24]}: the generalized hybrid family, its
loss model, and catalog-part selection;
(ii) \emph{extended}: the count optimization of [1], which sizes switch widths continuously for
one topology, now re-synthesizes the topology from discrete parts at every count; the
self-consistent electrothermal treatment of [18] is carried to an explicit mounted package with a
declared cooling boundary;
(iii) \emph{new}: the comparison of three placements of one synthesized regulator under identical
cooling with a feasibility filter; the thermoelastic solution with an all-site joint search, the
viscoplastic submodel and the optional interface feedback; and the per-net realizability check of
the tier order evaluated together with its thermal consequences. We also state what this work does
not do that [1] does (spatial IR drop, transient droop, automatic feedback into the synthesis).
Categorical statements (``no tool can'', ``invisible to any single-domain analysis'') were removed.
The ablation the reviewer suggested is now given explicitly in Section~VII: with electrical models
only (common temperature) Z2 leads Z1 by \GapIsoHotZtwo~W; adding the electrothermal fixed point
under one top plate reverses this order (Z1 ahead by \GapAZtwo~W) and removes Z3 (infeasible);
adding the mechanics confirms Z1 under the top plate, while under top and bottom plates the
electrically best integration (Z2) is the mechanically worst. \pending{update with coupled Z3}
\end{reply}

\begin{comment}
MC2 -- The framework is not fully bidirectionally coupled; classify each coupling direction,
identify neglected mechanisms, conditions; targeted sensitivity to contact resistance or interface
conductance.
\end{comment}
\begin{reply}
We agree that ``fully coupled'' was imprecise. Section~II now defines the three kinds of
interaction the chain uses and classifies every direction: loss and temperature are solved together
(bidirectional fixed point); temperature loads the mechanics and the critical joint drives the
local submodel one way; mechanical feedback is neglected in production, with the reason stated
for each path. The electrical feedback through thermal strain is two orders of magnitude below the
resistivity change already in the fixed point. The thermal feedback through interface geometry is
\emph{not} negligible in general, and we now quantify it: with the optional interface feedback
(Section~IV-F), the top TIM follows the solved warpage under a flat, rigid plate and the
electrothermal and thermoelastic solutions are iterated to convergence. Under cooling~A this bound
raises the loss of Z1 by \FbAZonedP~W and of Z2 by \FbAZtwodP~W, widening Z1's lead; under cooling~B
the effect is \FbBZonedP--\FbBZtwodP~W and changes no ranking. \pending{final feedback numbers}
Contact and intermetallic resistance are carried in the attach and pad terms as the all-solder
bound; damage, delamination and joint degradation are outside the intact-package scope, which is now
stated in Sections~II and~IX.
\end{reply}

\begin{comment}
MC3 -- The omitted Z3 redistribution losses can reverse a central conclusion.
\end{comment}
\begin{reply}
We agree; this was the most important point of the review. The land-side redistribution is now a
term of the coupled ledger: its loss $2\kappa I^2 R_s(T)$, with the shape factor $\kappa$ from the
two-dimensional sheet solution and the sheet resistance at the converged temperature, is deposited
in the substrate's land-side layer and iterated with the thermal solution, and the deformation, the
critical joint, $R^*$ and the voltage drop are recomputed. \pending{z3land: per-thickness loss,
drop, ranking, tie thickness, 30-mV thickness, board-plane bound}
\end{reply}

\begin{comment}
MC4 -- The Z1 design depends on input decoupling that is not in the physical stack.
\end{comment}
\begin{reply}
\pending{decap: requirement from the waveform, circuit simulation with ESR/ESL/DC bias and the
package network, parts, placement, k achieved, ESR loss, coupled result, Table~II row, Fig.~3
schematic}
\end{reply}

\begin{comment}
MC5 -- Z3 device temperatures exceed the published GaN limits; complete thermal-feasibility table.
\end{comment}
\begin{reply}
We agree. The framework now has an explicit feasibility stage (Section~II, Algorithm), and every
compared case is checked against: the GaN junction rating ($T_J \le 150\,^\circ$C for EPC2014C and
EPC2057), bounded from the tier maximum by the device's own $R_{\theta\mathrm{JC}}$ at its dissipated
power; the range of the on-resistance fit (25--125\,$^\circ$C); the X7R/X7S class of the capacitors
and the inductor rating (125\,$^\circ$C); the glass transition of the build-up dielectric and
encapsulant (about 150\,$^\circ$C) and of the board (about 170\,$^\circ$C); and the solder law's
range (125\,$^\circ$C) and the SAC305 solidus (Table~S5, Table~III). Z3 under the top plate alone
fails six of them (GaN junction \TjAZthree\,$^\circ$C) and is now classified as infeasible and
removed from the cooling-A ranking; Z2 under cooling~A passes with \TjMarginAZtwo~K of junction
margin but only 0.7~K below the capacitor class; every cooling-B case passes. Without local
decoupling (S1) the embedded design also fails (solder range, capacitor class). Fig.~8 marks the
150-$^\circ$C isotherm and Fig.~9 marks infeasible cases.
\end{reply}

\begin{comment}
MC6 -- The missing converter losses may not cancel between architectures.
\end{comment}
\begin{reply}
We agree that cancellation had to be shown. Dead-time conduction (EPC2057 reverse conduction over a
2--10-ns dead time, temperature dependent), inductor core loss (the material's loss line at the
actual flux ripple), AC winding resistance and capacitor dielectric loss were evaluated for every
compared case at its own tier temperatures and for every count of the screen (Section~VIII,
Table~S7). They add \OmitLo--\OmitHi~W (\OmitNom~W nominal) to every integration. They change the
Z1--Z2 gap under cooling~A by $+$\OmitGapAZoneZtwoLo{} to $+$\OmitGapAZoneZtwoHi~W and the Z2--Z3 gap
under cooling~B by \OmitGapBZtwoZthreeLo{} to \OmitGapBZtwoZthreeHi~W; no ranking flips under any
combination. The count optimum stays at $M = \OmitDseOpt$ (\OmitDseMarginLo--\OmitDseMarginHi~W
ahead) at 25 and 100\,$^\circ$C. \pending{authors: C$_{\mathrm{oss}}$ term and per-switch stress}
\end{reply}

\begin{comment}
MC7 -- The 76-cell optimum is not an ETM optimum; evaluate candidates with the complete chain;
report sampled counts; frequency grid; trade-off; constraints vs post-evaluation.
\end{comment}
\begin{reply}
We agree and now present the count study as an electrical screen followed by coupled evaluation
of the strongest candidates (Sections~III-C and~VIII). The \DseNcounts{} sampled counts are listed,
each jump of Fig.~4 is attributed to the constraint that becomes active, and the infeasibility
beyond $M = 76$ is explained (the area budget falls below the smallest rated device set). A
switching-frequency grid up to sixteen times finer changes no optimum and no count by more than
0.9~W. Temperature, manufacturability and mechanics are stated explicitly to be evaluated after
the synthesis, not constraints of it. \pending{countopt: best square candidate through the full
chain and the trade-off}
\end{reply}

\begin{comment}
MC8 -- Solder-joint assessment: parameters, cycles, cycle definition, constitutive accuracy,
inelastic strain is not life, benchmarking.
\end{comment}
\begin{reply}
\pending{solder: parameter table S6, 16-cycle stabilization, ramp times and justification of the
prescribed cycle, band sensitivity, alternative published SAC305 law, benchmark}
We now use ``response'' consistently and never ``reliability'' or ``life'' for $\deps$.
\end{reply}

\begin{comment}
MC9 -- Maximum displacement does not necessarily identify the joint with maximum inelastic response.
\end{comment}
\begin{reply}
\pending{solder: Pareto set in ($\Delta u$, $T_{\mathrm{j}}$) plus corners solved at L0 for every
compared case; site of largest response}
\end{reply}

\begin{comment}
MC10 -- Cooling model: physical and energy assessment.
\end{comment}
\begin{reply}
Section~IV-E and Table~S4 now give the complete specification: 223 channels of
$50\times300$~\textmu m and 22.36~mm length, water at 25\,$^\circ$C, 27~g\,s$^{-1}$ per plate ($Re =
\CoolRe$), the Shah--London correlation with its validity, fin efficiency and area amplification
defined in the rewritten Eq.~(3), and the entrance check: the thermal entrance length
(\CoolLth~mm) is comparable to the channel, so the fully developed Nusselt number is conservative,
consistent with the CFD comparison. The pressure drop is \CoolDpFd~kPa fully developed and about
\CoolDp~kPa with entrance and plenum losses; the pumping power is \CoolPhyd~W hydraulic per plate
(\CoolPshaftLo--\CoolPshaftHi~W at the shaft). Under cooling~B each plate receives the full flow,
so the comparison is one of identical hardware; pumping power is charged in the new
platform-inclusive accounting, which shows that the second plate costs about \PumpShaftPlate~W and
is a mechanical rather than an efficiency decision for Z1. The heat balance per plate is reported
(Table~II). A non-uniform die power map (75\,\% of the power in the central half of the die)
raises the die peak by 15--20~K but changes the losses by less than 0.3~W and no ranking
\pending{hotspot mechanics}. While documenting the boundary we found that the plate coefficient
includes the caloric resistance of a uniform sink while the solver also resolves the coolant warming;
this counts half the coolant rise (\CoolCalK~K per kW) twice. Removing it lowers every loss by
1--5~W and changes no ranking (Section~VII); the production results are kept as the conservative
case and the correction is reported.
\end{reply}

\begin{comment}
MC11 -- Feasibility checks do not cover the full interconnect and assembly.
\end{comment}
\begin{reply}
Section~IV-D now states the declared rule set (via formation in a build-up dielectric laminated
around the components, capture planes and pads, overlay allowance, the 1:1 aspect-ratio limit),
the condition under which the IPC-2226A practice transfers to the inter-tier dielectric, and
distinguishes geometric feasibility under declared rules (what is shown), process feasibility
(not shown) and assembly feasibility (not evaluated); the same distinction is made for the
30-\textmu m-pitch Cu-pillar attach and the board opening. The ball-provisioning problem of Z3 is now
a feasibility failure in its own right: 1334 of 1369 sites leave 35 for every signal and auxiliary
rail, so Z3 is infeasible in this package irrespective of the electrical comparison unless the
ball pitch falls to \PitchMaxZthree~mm (Section~VII, Table~III). \pending{flying-capacitor tier}
\end{reply}

\begin{comment}
MC12 -- Verification is not sufficient for full-system accuracy; complete-package cross-check.
\end{comment}
\begin{reply}
We agree and have done both. Section~V now opens with Table~IV, which states for each kind of
check what it establishes and what it does not, and Fig.~10 shows only matched-reference checks.
A complete-package cross-check was added: the full Z1 package (all 16 layers, nine materials,
\FullPkgElemsZone{} hexahedra, the converged heat map and a uniform plate film) solved by the
production path and by MOOSE on the identical mesh agrees to \FullPkgErrZone~K (Z3:
\FullPkgErrZthree~K). MOOSE solutions on refined copies of up to \FullPkgMaxElemsZone~M hexahedra
showed that the production mesh leaves the converter stack and board of Z1 2--3~K cold (the die is
converged), converging at order 0.6; our discretization uncertainty was revised upward
accordingly (MC13).
\end{reply}

\begin{comment}
MC13 -- The architecture ranking needs a more rigorous uncertainty treatment.
\end{comment}
\begin{reply}
The new Table~V lists $U$ for every compared output of every architecture in both cooling tiers,
built from the extrapolated discretization error (at the order observed in the MOOSE ladder), the
through-thickness split, the fixed-point tolerance and the matched-reference agreement, and states
the rule: a gap is resolved when it exceeds $2(U_a+U_b)$, and resolved with correlation credit when
it exceeds twice the paired estimate in which the signed discretization errors of two architectures
on one mesh family cancel. The 3.1-W Z1--Z2 gap under cooling~A is resolved only with correlation
credit (the gap changes by 0.04~W from r21 to r31). The reviewer was right about the solder
response: under cooling~B the Z1--Z3 difference is not resolved, and neither is their warpage
difference; the text now says so. Model-form uncertainty (material cards, boundaries, omitted
physics) is handled by rerunning both architectures under each declared alternative (Fig.~11),
and a resolved ranking that an alternative reverses is reported as conditional.
\pending{update with the solder study's stabilized values}
\end{reply}

\begin{comment}
MC14 -- The DC rail-aware ledger does not characterize high-frequency power integrity.
\end{comment}
\begin{reply}
We agree and restrict the claims explicitly to steady-state conduction loss, temperature,
deformation and joint response of an intact package (Sections~I and~IX). The ledger now states, net
by net, which current it charges: the RMS currents of the converter model in the switched nets, the
scenario RMS $kI$ in the 48-V nets wherever the pulsed waveform flows, and the RMS of the cell output
current in the 1-V path. Frequency-dependent resistance, high-frequency return paths, droop and
decoupling resonance are outside the ledger; the input network's resonance and impedance are
checked by the circuit simulation of the decoupling design (MC4). \pending{decap circuit results}
\end{reply}

\begin{comment}
MC15 -- Key geometry, material and solver parameters are missing.
\end{comment}
\begin{reply}
The new supplementary material lists every material card with its value and status, every layer of
every stack as meshed, mesh sizes and degrees of freedom, solver settings and tolerances, the
hardware and software, runtimes by stage, the cold plate, the feasibility limits, the solder model,
every varied input with its range, the S0--S2 states of the selected regulator, the verification
matrix and the uncertainty budget. All tables are generated from the solver inputs and outputs.
\end{reply}

\begin{comment}
MC16 -- System boundaries and external power-feed assumptions.
\end{comment}
\begin{reply}
Eq.~(1) is now fully defined, including where the RMS factor $k$ applies: wherever the pulsed
current flows, in the package 48-V network of Z1 and in the external feed of Z2 and Z3 alike. Two
physical feeds are evaluated: for Z2, 48~V through the BGA ring and a substrate 48-V RDL to the
converter, solved by the same voltage-fed network as Z1 (loop \FeedRbgaZtwo~m$\Omega$,
\FeedPbgaZtwo~W at S0, \FeedPbgaZtwoSone~W without local decoupling); and for Z3, the board's own
48-V plane pair (\FeedRZthree~m$\Omega$). Both lie far below the cooling-B crossovers
(\RstarTierBZtwo{} and \RstarTierBZthree~m$\Omega$). Platform-inclusive losses including the feed and
the pumping power of every plate are reported for both tiers (Table~III). \pending{update after Z3}
\end{reply}

\section*{Figures}
\begin{itemize}
\item \textbf{Fig.~1:} coupling directions distinguished (double, single, dotted arrows), the
package-internal boundary and the external feed marked, text enlarged. \pending{figure}
\item \textbf{Fig.~2 (converter):} local input capacitor added; input-current waveform with
amplitude and RMS from the circuit simulation. \pending{decap}
\item \textbf{Fig.~3 (DSE):} sampled counts marked, each jump labeled with its cause, the
infeasible range shaded, coupled-chain candidates marked; screen and package baselines
reconciled in the text (96.1 against 94.1~W: same conversion loss, different routing model).
\pending{figure}
\item \textbf{Fig.~4 (sections):} input decoupling, return paths and critical dimensions; an
enlarged inter-tier connection. \pending{figure}
\item \textbf{Fig.~5 (realizability):} generation of the 144 assignments and the via construction
explained in Section~IV-D. \pending{inset}
\item \textbf{Fig.~6 (verification):} only matched-reference checks; mesh and solder-family results
moved to Table~V.
\item \textbf{Fig.~7 (fields):} full temperature scale, the 150-$^\circ$C device rating marked,
critical joints enlarged. \pending{figure}
\item \textbf{Fig.~8 (loss):} complete and bounded terms distinguished; waveform assumptions in the
caption. \pending{figure after Z3}
\item \textbf{Fig.~9 (comparison):} panel (e) described in the caption, L0 defined, $\deps$ stated to
be a relative response, infeasible cases marked, $U$ shown. \pending{figure}
\item \textbf{Fig.~10 (robustness):} the criterion stated numerically and the range of each gap
under each decision added. \pending{figure}
\end{itemize}

\section*{Tables}
\textbf{Table~I (specification)} now defines $J_{\mathrm{t}}$ as a minimum current density (the cell
uses 5.61 of its \DseFortyNineBudget~mm$^2$, i.e., \DseFortyNineJ~A\,mm$^{-2}$), gives each device's
rating against its operating stress, the inductor's saturation current against its peak current,
the flying capacitors' voltages, the cooling specification and the input decoupling.
\pending{flying-capacitor part and decoupling row}
\textbf{Table~II (comparison)} is reorganized into calculated performance, completeness of the loss
accounting and feasibility, with the maximum junction temperature, the bounded or coupled terms,
the voltage drop, the feed and platform-inclusive loss, the uncertainty of the joint response and
the heat balance per plate. \pending{table}

\section*{Equations}
Eq.~(1): $P_{\mathrm{int}}$ defined, $k$ justified and applied wherever the waveform flows, the
thermal treatment of the feed loss stated, two physical feeds evaluated. Eq.~(2): $J_{\mathrm{d}}$
(100~A\,mm$^{-2}$), the liner and the copper radius defined, contact resistance stated, the current
type of every net given. Eq.~(3): rewritten exactly as implemented, with the correlation, its validity,
$\eta_{\mathrm{f}}$, $\Lambda$, the caloric and base resistances, fluid properties, entrance effects
and hydraulics. Eq.~(4): sampling and reference state stated; validated against the local response
at competing sites (MC9). Eq.~(5): \pending{solder: law, calibration, time stepping, band meaning and
sensitivity, stabilization}

\section*{Numerical audit}
\begin{itemize}
\item 96.1 vs 94.1~W: the screen and the package ledger share the conversion loss
(\DseFortyNineConv~W) and differ only in the routing model; now stated in Section~III-C.
\item 88.4 vs 88.5~W: the value is \DseOptLoss~W in both the figure and the text (generated).
\item $J$ = 3.64 vs 3.5~A\,mm$^{-2}$: $J_{\mathrm{t}}$ is a minimum; defined in Section~III-B and Table~I.
\item 5.71 vs 5.80$\times10^{-3}$: 5.71 is the finest local mesh (L3) of the evaluation and 5.80 the
L0 level used for every comparison; both are now labeled. \pending{solder update}
\item Z3 at 179.7\,$^\circ$C: now classified infeasible (MC5).
\item Z3's 18.1-W advantage: re-evaluated with the coupled land side (MC3). \pending{z3land}
\item 1334 of 1369 balls: a feasibility failure (MC11).
\item Mesh criterion vs the 3.1-W gap: per-output budget and rule (MC13).
\item Fig.~9 panel (e): described.
\item S2: ``infeasible under the intact-solid solder model'' (Section~VI).
\end{itemize}

\section*{Minor comments}
\begin{enumerate}
\item Abbreviations and symbols defined at first use (L0, $J_{\mathrm{d}}$, $\Lambda$, $U$, Z1--Z3).
\item ``Package-backside'' (Z2) and ``board-backside'' (Z3) used throughout; ``substrate backside''
appears only as a physical description.
\item Thermal, electrical and feed resistance distinguished by symbol and adjective.
\item Efficiencies labeled by boundary (package-internal or platform-inclusive).
\item The GaN coefficients are fitted over 25--125\,$^\circ$C; extrapolation beyond is flagged by the
feasibility stage.
\item The 1.5-A-per-ball rule is a derated continuous design target, stated as such.
\item Legal sites and power/return assignment per architecture: Table~S3. \pending{table}
\item Encapsulant properties: Table~S1 ($k = 0.7$~W\,m$^{-1}$K$^{-1}$, $E = 8$~GPa, 30~ppm/K).
\item The board's copper is homogenized by a verified mixing rule; stiffness is
temperature-independent (Table~S1).
\item The evaluation time includes the fixed point, the structural solve with the joint search and
the local submodel; geometry and meshing take under two minutes; sensitivities are extra (Section~V).
\item The fixed-point tolerance enters $U$ explicitly (Table~V).
\item The range of every varied input: Table~S7.
\item S0--S2 states of the selected regulator: Table~S8.
\item ``Reliability'' is no longer used for $\deps$.
\item Legends, notation and caption coverage checked throughout.
\end{enumerate}
\end{document}
